% Draft: rationale of the evaluation design (fixed held-out test, nested
% selection, seeds, paired bootstrap, cluster-level tests) versus repeated
% cross-validation schemes such as 5x2cv.
% Target: Chapter 7 (Experimental Protocol), replication and inference sections.
% Status: draft v1, 2026-11-02. References checked against Crossref DOIs.
% Project sources: specs/statistical_evaluation_protocol.md,
%   configs/evaluation_protocol_v3.yaml, specs/homology_grouping_protocol.md.

\section{Rationale of the Evaluation Design}\label{draft:evaluation-design}

\draftmeta{Chapter~7 (replication and inference)}
{Draft v1 --- describes the active protocol; group-aware folds await evaluation protocol v4}
{\texttt{specs/statistical\_evaluation\_protocol.md};
\texttt{specs/homology\_grouping\_protocol.md}}

\subsection{Design}

Each dataset is evaluated on a held-out test partition that is used once,
after the model configuration has been locked. Within the training partition,
a fixed five-fold stratified manifest is used only for hyperparameter and
epoch selection. Stochastic models are trained with ten predeclared seeds;
deterministic models are fitted once. The primary outcome is the Matthews
correlation coefficient (MCC), which uses every cell of the confusion matrix
and remains informative under class imbalance \citep{chicco2020mcc}. Within a
dataset, models are compared by the median of their seed-paired MCC
differences, with a 95\% interval from a paired bootstrap of the test records.
Across datasets, each dependency cluster contributes one effect, and models are
compared with the Wilcoxon signed-rank test, or the Friedman test followed by
predeclared pairwise comparisons, with Holm correction
\citep{demsar2006statistical}.

\subsection{Why repeated cross-validation tests are not used}

Repeated cross-validation schemes produce several performance estimates per
dataset, but these estimates are not independent: training sets overlap
between folds and between repetitions. Treating them as independent samples in
a paired $t$-test underestimates the variance and inflates the type~I error
\citep{dietterich1998approximate,nadeau2003inference}. Two families of
corrections exist. The $5\times2$ cross-validation paired $t$-test
\citep{dietterich1998approximate} and its combined $F$ variant
\citep{alpaydin1999combined} use five repetitions of two-fold cross-validation
and were calibrated for that exact design. The corrected resampled $t$-test
inflates the variance estimate to account for the overlap between training sets
\citep{nadeau2003inference}.

The protocol uses neither, for four reasons that hold for any dataset it
admits.
\begin{enumerate}
  \item \textbf{Comparability with published benchmarks.} Fixed released
    partitions are preserved so that results remain comparable with the
    literature. Re-partitioning every dataset into repeated halves would
    discard that comparability.
  \item \textbf{Training-set size.} Two-fold schemes train on half of the
    data. When performance still improves with sample size, the estimate is
    pessimistically biased \citep{efron1997improvements}. Because learning
    curves differ between model families, the bias need not be equal for
    capsule networks and linear baselines, and it could therefore distort
    their comparison rather than only shift both.
  \item \textbf{Separation of selection and evaluation.} Selecting
    hyperparameters on the same folds used to estimate performance biases the
    estimate optimistically \citep{varma2006bias}. Keeping selection inside the
    training partition and evaluating once on the held-out partition removes
    that bias without the cost of a full nested repetition.
  \item \textbf{Unit of inference.} The claim of interest is whether one model
    family outperforms another across promoter datasets. Statistical tests for
    comparing classifiers over multiple datasets treat datasets as the sample
    units \citep{demsar2006statistical}; the protocol strengthens this by
    collapsing related datasets into one dependency cluster.
\end{enumerate}

\subsection{Roles of folds, seeds, and the bootstrap}

The three sources of repetition answer different questions, and none of them
is promoted to an independent observation.
\begin{itemize}
  \item \textbf{Inner folds} serve only model and epoch selection.
  \item \textbf{Seeds} measure algorithmic instability (initialisation, sampling
    and training order) conditional on one dataset. They are not biological
    replicates, so seed-level differences are summarised by their median, not
    tested as independent samples.
  \item \textbf{The bootstrap} measures sampling uncertainty of the test set.
    The same resampled indices are applied to every model, so that intervals
    describe paired differences. When records are dependent, for example
    through homologous or overlapping windows, records are resampled by group
    rather than individually \citep{field2007bootstrapping}, and partitions and
    inner folds are built so that no group is split
    \citep{teufel2023graphpart,petti2022benchmark}.
\end{itemize}

\subsection{Limitations}

A single held-out partition gives no estimate of variation due to the choice of
partition itself; the bootstrap captures test-set sampling variation only.
Seed variation captures part of the training variation but not that of a
different training sample. Both limitations are stated alongside every
interval, and the cross-dataset test, not any within-dataset interval, carries
the confirmatory claim.
